\documentclass[11pt]{article}
\usepackage{graphicx} % Required for inserting images
\setlength{\parindent}{0pt}
\usepackage[bookmarks=true, colorlinks=true, urlcolor=blue, linkcolor=black]{hyperref}

\usepackage{enumitem}
\usepackage[utf8]{inputenc}
\usepackage[T1]{fontenc}
\usepackage[english]{babel}
\usepackage{geometry}
\usepackage{lastpage}
\usepackage{fancyhdr}
\usepackage{datetime}
\usepackage{parskip}
\usepackage[normalem]{ulem}
\usepackage{xurl}
\urlstyle{same}


\usepackage[absolute]{textpos}
\usepackage{soul}
\usepackage{tikz}
\usetikzlibrary{calc}

% Link to other SOPs, forms
\usepackage{xr}
\externaldocument[102-]{../2_medical/sop_102_medical}
% Definitions for SOP sections 3.
\newcommand{\AdverseEvent}{\textbf{Adverse Event}: Any untoward or unfavorable medical occurrence in a human subject, including any abnormal sign, symptom, or disease, temporally associated with the subject's participation in the research, whether or not considered related to the subject's participation in the research.}

\newcommand{\ControlRoom}{\textbf{Control Room}: B17; The main room of the MRI suite that contains the computer displays and keyboards that control the MRI scanner. The Control Room is accessed from B15 that is controlled by an Electronic Lock.}

\newcommand{\DcmComputer}{\textbf{DICOM Computer}: The Mac Studio located in the Control Room the receives DICOM files from the scanner.}

\newcommand{\ElectronicLock}{\textbf{Electronic Lock}: All doors at the ND-HNC are controlled by electronic locks. The electronic lock system recognizes the owner's Irish1Card, logs the time of the attempted entry, and unlocks the door if the owner has been granted access.}

\newcommand{\EEG}{\textbf{EEG}: Electroencephalograpy, a technique used to measure changes in electromagnetic properties of neural tissue.}

\newcommand{\EquipmentRoom}{\textbf{Equipment Room}: B18; The room containing hardware for MRI operations.}

\newcommand{\FDA}{\textbf{FDA}: Food and Drug Administration.}

\newcommand{\GoodCatch}{\textbf{Good Catch}: Early recognition of an event with the potential to cause damage, injury, or illness.}

\newcommand{\Helper}{\textbf{Helper}: Second person in Zones 3 and 4 who assists in emergency responses, typically Level 2-3.}

\newcommand{\Incident}{\textbf{Incident}: An occurence, action, or event that warrants reporting (e.g. a policy is not followed and/or an individual's health is at risk). Includes Adverse Events, Near Misses, Unanticipated Problems, etc.}

\newcommand{\Lead}{\textbf{Lead}: Primary person in Zones 3 and 4 responsible for the Visitor's safety, in charge during emergency responses. Level 4 operator or MRI Technologist.}

\newcommand{\MagnetRoom}{\textbf{Magnet Room}: B17A; The room within the MRI suite that contains the MRI scanner. It is accessed from within the MRI Control room (B17) through a doorway that is controlled by an Electronic Lock.}

\newcommand{\NearMiss}{\textbf{Near Miss}: An incident that occurred but did not cause personal injury/illness e.g. a small/contained fire, trip on sidewalk or carpeting, falling object.}

\newcommand{\NDHNC}{\textbf{ND-HNC}: Notre Dame Human Neuroimaging Center, Lower Level of Veldman Family Psychology Clinic.}

\newcommand{\Operator}{\textbf{Operator}: MRI Technologist or Level 4 researcher, responsible for conducting research operations and maintaining equpiment. The Lead during emergency responses.}

\newcommand{\PIN}{\textbf{PIN}: Personal Identification Number, separate for Irish1Card and security.}

\newcommand{\ResearchAssistant}{\textbf{Research Assistant}: Individual with access Levels 1-3 to the MRI research suite. Typically the Helper during emergency responses.}

\newcommand{\SeriousInjury}{\textbf{Serious Injury}: An injury or illness that is life threatening; results in permanent impairment of a body function or permanent damage to a body structure; or necessitates medical or surgical intervention to preclude permanent damage or impairment.}

\newcommand{\StimControl}{\textbf{Stimulus Control Computer}: The Windows machines in the Control Room that deploy fMRI tasks and capture participant responses.}

% \newcommand{\Subject}{\textbf{Subject}: An individual who is participating in an experimental protocol at ND-HNC and has not taken the ND-HNC safety course.}

\newcommand{\Supervision}{\textbf{Supervision}: A supervisor is continuously present in both sight and sound range of the individual or task being monitored. This approach ensures the supervisor can actively observe, direct, and immediately intervene if a complication or safety issue arises.}

\newcommand{\TMS}{\textbf{TMS}: Transcranial magnetic stimulation, a technique used to disrupt normal operations of neural tissue.}

\newcommand{\TornadoWarn}{\textbf{Tornado Warning}: A tornado has been sighted or indicated by weather radar and there is imminent danger to life and property. Action is required.}

\newcommand{\TornadoWatch}{\textbf{Tornado Watch}: Tornados are possible in and near the watch area. Be ready to act quickly and move to the designated tornado safe area.}

\newcommand{\UnanticipatedProblem}{\textbf{Unanticipated Problem}: An incident, experience, or outcome that meets \textbf{all} of the following criteria: (1) unexpected (in terms of nature, severity, or frequency) given (a) the research procedures that are described in the protocol-related documents, such as the IRB-approved research protocol and informed consent document; and (b) the characteristics of the subject population being studied; (2) related or possibly related to participation in the research; and (3) suggests that the research places subjects or others at a greater risk of harm than was previously known or recognized.}

\newcommand{\VFPC}{\textbf{VFPC}: Veldman Family Psychology Clinic.}

\newcommand{\Visitor}{\textbf{Visitor}: An individual who has not taken the ND-HNC MRI Safety Course but has a legitimate reason for entering the MRI suites. Research participants and their legal guardians are Visitors. Other Visitors may include: individuals providing/receiving a tour of the VFPC and ND-HNC, faculty and staff who have not received MRI Safety Course, and research assistants.}


\usepackage[super]{cite}
\hypersetup{colorlinks=true, citecolor=black}

% figures and tables
\usepackage{graphicx}
\graphicspath{{../../../../media/sops}{../../../../media/nd_brand}}
\usepackage{rotating}
\usepackage{float}
\usepackage{subfig}
\usepackage{caption}
\DeclareCaptionLabelFormat{adja-page}{\hrulefill\\#1 #2 \emph{(previous page)}}

% Branding
\providecommand\fromlogo{nd_logo.png}
\definecolor{NDblue}{RGB}{12,35,64}
\definecolor{NDgold}{RGB}{174,145,66}
\definecolor{NDgray}{RGB}{193,205,221}

% Center info
\providecommand\nameuniversity{University of Notre Dame}
\providecommand\namecenter{Human Neuroimaging Center}
\providecommand\addrcenter{501 N Hill St. South Bend, IN. 46617}
\providecommand\namecollege{College of Arts and Letters}
\providecommand\addrweb{https://neuroimaging.nd.edu}

% \renewcommand{\headrulewidth}{0pt}
% % \newdateformat{yearmonth}{\THEYEAR-\THEMONTH-\THEDAY}
% \pagestyle{fancy}
% \fancyhead[L]{}
% \fancyhead[R]{Approved: TBD}
% \fancyfoot[C]{}
% \fancyfoot[L]{ND-HNC SOP 103: v.0.1.0}
% \fancyfoot[R]{p. \thepage /\pageref*{LastPage}}


% Format header
\urlstyle{sf}
\pagestyle{fancy}

\fancyhead{}
\fancyhead[L]{%
    % LOGO
    \begin{textblock*}{2in}[0.3066,0.39](1.0in,0.5in)
        \includegraphics[width=2in]{\fromlogo}
    \end{textblock*}

    % Above line
    \begin{textblock*}{6.375in}(1.5in,0.32in)
        \sffamily
        \hfill \textcolor{NDgold} \namecenter\ | \textcolor{NDgold}\namecollege
    \end{textblock*}

    % Below line
    \begin{textblock*}{6.375in}(1.5in,0.53in)
        \sffamily
        \hfill \textcolor{NDgold} \addrcenter\ | \textcolor{NDgold}\addrweb
    \end{textblock*}

    %Line
    \begin{tikzpicture}[remember picture,overlay]
        \draw[color=NDblue,line width=0.7pt] (current page.north west)+(2.45in,-0.48in) -- ($(-0.625in,-0.48in)+(current page.north east)$);
    \end{tikzpicture}
}
\renewcommand{\headrulewidth}{0pt}

% Format footer
\fancyfoot[L]{v.0.2.0}
\fancyfoot[C]{\thepage}
\fancyfoot[R]{\textbf{ND-HNC SOP: 103}}

\renewcommand{\footrulewidth}{0.5pt}



% start document
\begin{document}

\renewcommand{\labelenumii}{\arabic{enumi}.\arabic{enumii}}
\renewcommand{\labelenumiii}{\arabic{enumi}.\arabic{enumii}.\arabic{enumiii}}
\renewcommand{\labelenumiv}{\arabic{enumi}.\arabic{enumii}.\arabic{enumiii}.\arabic{enumiv}}

\begin{center}
	\textbf{\Large SOP 103: Reporting Incidents}\\
	\hrulefill
\end{center}


\section{Summary}
This document establishes procedures for Incident and Adverse Event reporting at the Notre Dame Human Neuroimaging Center (ND-HNC). It is intended to match the Incident or Adverse Event with the entity to which reports are made.


\section{Scope}
These policies apply to all investigators, research assistants, and staff who use the ND-HNC.

\section{Definitions}

\begin{itemize}
	\item \AdverseEvent\cite{ohrp2007}
	\item \GoodCatch
	\item \Incident
	\item \NearMiss
	\item \SeriousInjury\cite{injury}
	\item \UnanticipatedProblem\cite{ohrp2007}
	\item \Visitor
\end{itemize}

\section{Notes}

This SOP does not replace any requirements or obligations research personnel have for reporting Adverse Events or Unanticipated Problems.

\section{Policies \& Procedures}

The safety of research personnel and participants is of paramount importance in all research activities. Accordingly, any Incident that puts an individual's health, safety, or well-being at risk must be reported and documented promptly. Center (ND-HNC), University, federal, and manufacturer policies require different levels of reporting at different time intervals; a single Incident is likely to be reported to multiple entities.

An \uline{Adverse Event} can be understood as any accident, injury, illness, or health-related complaint that may be \textit{\textbf{temporally}} associated with a Visitor's involvement with the ND-HNC. These may range from minor (e.g. headache) to serious (see, Serious Injury). For instance, a Visitor may have a tour of the center and then complain about feeling faint. We do not differentiate between expected and unexpected Adverse Events.

A \uline{Serious Injury} describes a severe medical emergency. It includes any injury or illness that is immediately life-threatening, causes lasting damage to how a person's body functions or looks, or requires urgent medical or surgical treatment to prevent that permanent damage from happening.

An \uline{Unanticipated Problem} is an unexpected event or outcome that is related to research operations and constitutes a new, previously unknown risk. The emergence of Lululemon's Silverscent material, and the resulting burns, are an example.

Because the distinction among the various Incident types may not always be immediately clear, personnel should prioritize safety and appropriate response over classification. Any Incident must be addressed promptly, documented appropriately, and reported to the relevant University authorities as required. Final determination of Incident classification may be made following review of the circumstances.

The ND-HNC follows all guidelines and reporting measures set in place by the University of Notre Dame's Incident Reporting and Management Program from the Risk Management and Safety Department\cite{riskman}.

%
\subsection{Incidents involving Faculty, Staff, Students \& Contractors}
\label{ssec:policy-staff}
Any Incident which occurs with an employee or student of the University of Notre Dame must be immediately reported to ND-HNC staff.

ND-HNC staff then:
\begin{enumerate}
	\item Navigate to Risk Management\cite{riskman}
	\item Follow the prompts, selecting the appropriate option (e.g. Good Catch)
	\begin{enumerate}
		\item If reporting an injury, select \textit{\textbf{``Yes''}} to complete the \textbf{\underline{Safety Incident Report}}.
		\item Note: ND faculty or staff must submit this report on behalf of unpaid students and contractors.
	\end{enumerate}
	\item Log the Incident in ND-HNC Form LG03: Incidents.
\end{enumerate}

%
\subsection{Incidents involving Visitors}
\label{ssec:policy-visitor}
Any Incident which occurs with an individual not affiliated with Notre Dame, or with an affiliate during their personal time (e.g. lunch break) must be immediately reported to ND-HNC staff.

ND-HNC staff then:
\begin{enumerate}
	\item Navigate to Risk Management\cite{riskman}
	\item Follow the prompts, selecting the appropriate option (e.g. Good Catch)
	\begin{enumerate}
		\item If reporting an injury, select \textit{\textbf{``No''}} to complete the \textbf{\uline{General Liability Incident Report}}.
	\end{enumerate}
	\item Log the Incident in  ND-HNC Form LG03: Incidents.
\end{enumerate}

%
\subsection{MRI-Related Incidents}
Incidents related to the MRI scanner are reported to: (1) the University, (2) the Food and Drug Administration (FDA), and (3) the manufacturer (Siemens)\cite{medrep}.

\subsubsection{University Reporting}
Follow procedures detailed in Sections \ref{ssec:policy-staff}, \ref{ssec:policy-visitor}.

\subsubsection{FDA Reporting}
ND-HNC staff voluntarily report Adverse Events and Serious Injuries to the FDA through their MedWatch\cite{medwatch} program. Items that should be reported to the FDA:

\begin{itemize}
	\item Device-related Death: within 10 work days of becoming aware (Form FDA 3500A).
	\item Annual summary of death and serious injury reports: January 1 for the preceding year (Form FDA 3419).
	\item Unexpected side effects or adverse events including but not limited to, skin rashes to more serious complications.
	\item Product quality problems such as information if a product isn't working properly or if it has a defect.
	\item ``Information that reasonably suggests'' that a medical device has caused or contributed to a reportable event including information such as professional, scientific or medical facts and observations or opinions.
	\item ``Information that is reasonably known'' to user facilities including information found in documents and any information that becomes available as a result of reasonable follow-up within the facility.
\end{itemize}

To file a report:

\begin{enumerate}
	\item Navigate to MedWatch\cite{medwatch}.
	\item Select `Report a Problem'.
	\item Select `Health Professional'.
	\item Fill out the  \uline{\textbf{MedWatch Voluntary Health Professional Report}} to the best of your ability.
\end{enumerate}


\subsubsection{Siemens Reporting}
ND-HNC staff voluntarily report Adverse Events and Serious Injuries to the MRI Manufacturer (Siemens). Items that should be reported to Siemens:

\begin{itemize}
	\item Device-related Death: within 10 work days of becoming aware (Form FDA 3500A).
	\item Device-related Serious injury: within 10 work days of becoming aware (Form FDA 3500A).
\end{itemize}

To file a report:

\begin{enumerate}
	\item Call Siemens Customer Care Center: 1-800-743-6367\cite{siemens-care}.
	\item Alternatively, call Siemens Healthineers: 1-800-888-7436\cite{siemens-health}.
\end{enumerate}

\section{Figures}

None.

\begin{thebibliography}{999}
	\bibitem{ohrp2007}
		Office for Human Research Protections. (2007).
		Reviewing and reporting unanticipated problems involving risks to subjects or others and adverse events: OHRP guidance (2007).
	\bibitem{injury}
		21 U.S.C. § 802(25).
	\bibitem{riskman}
		\url{https://riskmanagement.nd.edu/incident-reporting-management}
	\bibitem{medrep}
		\url{https://www.fda.gov/medical-devices/medical-device-safety/medical-device-reporting-mdr-how-report-medical-device-problems}
	\bibitem{medwatch}
		\url{https://www.fda.gov/safety/medwatch-fda-safety-information-and-adverse-event-reporting-program}
	\bibitem{siemens-care}
		\url{https://www.siemens.com/en-us/contact/}
	\bibitem{siemens-health}
		\url{https://www.siemens-healthineers.com/en-us/how-can-we-help-you}
\end{thebibliography}

\end{document}